\noindent\textbf{(a)} Sean \(\mu_n\in((n-1)\pi,(n-1)\pi+\pi/2)\) las raíces de \(\mu_n\tan\mu_n=1\), n\ensuremath{\ge}1.
\[\lambda_n=\mu_n^2,\qquad \phi_n(x)=d_n\cos(\mu_nx),\quad d_n=\sqrt{\frac2{1+\sin^2\mu_n}}.\]

\noindent\textbf{(b)} \[1-x=\sum_{n=1}^{\infty}c_n\phi_n(x),\qquad
c_n=\frac{d_n(1-\cos\mu_n)}{\mu_n^2}.\]

\noindent\textbf{(c)} \[f\sim\sum_{n=1}^{\infty}c_n\phi_n,\qquad
c_n=\frac{d_n\sin(\mu_n/2)}{\mu_n}.\]
La suma vale1/2 en x=1/2; en los puntos de continuidad coincide con f.
